\begin{reviewcomment}{7}
	Can you expand on the real cost of randomizing gas padding?
	The gas padding ablation is important, as it shows that policy-dependent gas signals are more recoverable when the padding is removed.
	But it costs about 38\% in access-gas and 40\% in commit-gas overhead to provide this protection.
	The manuscript would benefit from a more explicit discussion of this privacy-cost trade-off, e.g.\ whether the overhead remains practical for realistic transaction volumes and whether the padding strategy remains effective when policies or circuits become more complex.
\end{reviewcomment}

\begin{reviewresponse}{7}
	We expanded the cost discussion in three ways.

	\textbf{Re-measured overhead.}
	The ablation was re-run on the revised system and the new corpus (a $1{,}119$-execution subset, executed without padding three days after the padded of-record run, on a fresh deployment).
	The overhead is now $46.1\%$ on access gas ($723{,}462 \to 1{,}056{,}674$) and $38.4\%$ on commit gas ($219{,}737 \to 304{,}098$).
	The increase on access relative to the submitted version has two causes that the manuscript now explains: the revised relation adds an on-chain Poseidon recomputation of the context digest, and Polygon Amoy prices the pairing precompiles higher than the local EVM (the same verifier costs $243{,}000$ gas locally and $372{,}000$ on Amoy), so the padding band had to be positioned higher.

	\textbf{Practicality at volume.}
	The overhead is a fixed premium per transaction---roughly $330{,}000$ gas on access and $85{,}000$ on commit at the band used here---so it scales linearly with transaction volume and does not compound.
	At the $30$~gwei tip of the evaluation a padded access transaction costs about $0.032$~POL against $0.022$~POL unpadded.
	The Discussion now states this together with the deployment rule the band follows: its base must lie above the bare access cost measured on the target network, which is calibrated once per network, and its width ($\pm 10\%$ here) determines how large a deterministic difference it can submerge.

	\textbf{Effectiveness under more complex policies or circuits.}
	Neither the bare cost nor the band depends on the policy: the fixed circuit gives every proof the same verification work and the same five public signals, so a more complex policy changes nothing on-chain.
	To show how the prover-side cost scales, we added an offline sweep of the program length ($L = 32$, $64$, $128$ with $N$ and $M$ fixed; new cost-scaling table): constraints grow from $65{,}000$ to $134{,}000$ and proving time from $0.38$ to $0.73$~s per proof on a laptop-class prover, while the verifier contract, its gas, and the padding band are unchanged.
	Richer predicates would enter through the predicate slots and their gadgets, again adding prover work and nothing on-chain, because predicate parameters are bound through the fixed-format parameter commitment rather than exposed as public inputs.
	The revised ablation section also identifies what the padding must absorb: the un-padded channel is a sub-kilogas calldata-pricing difference in a public argument, which a $\pm 105{,}000$-gas band submerges regardless of its origin; padding cannot hide a difference larger than its band, and we say so.

	\RevisionLocation{Section ``Discussion,'' new paragraph ``The cost of padding''; Section ``Privacy Evaluation,'' subsection ``Ablation,'' paragraph ``Cost,'' and the new cost-scaling table.}
	\begin{revisiontext}
		Three observations bound its practical weight.
		First, the overhead is a fixed premium per transaction---roughly $330{,}000$ gas on access and $85{,}000$ on commit at the band used here---so it scales linearly with transaction volume and does not compound: at the $30$~gwei tip of the evaluation, a padded access transaction costs about $0.032$~POL against $0.022$~POL unpadded.
		Second, neither the bare cost nor the band depends on the policy: the fixed circuit gives every proof the same verification work and the same five public signals, and the prover-side cost of longer or richer policies is paid off-chain.
		Third, the band must be positioned above the bare cost on the target network, which differs between networks [\ldots]; a deployment therefore calibrates the base target once per network, and the width of the band, $\pm 10\%$ here, is what determines how large a deterministic difference it can submerge.
	\end{revisiontext}
\end{reviewresponse}
